% !LW recipe = make-deck
% Level 4F (programs track, FPIL): comparing rows. Theory: equality tests,
% membership, difference, joins. Lean: dedup, diff, inter, joinOn.
% Build: make (rules in ../common/deck.mk).
\DocumentMetadata{
  lang        = en-GB,
  pdfstandard = ua-2,
  tagging     = on,
}
\documentclass{ltx-talk}

\input{preamble}

\title{Level 4F: Difference and Joins, by Comparing Rows}
\author{The Proofs and Programs Club}
\institute{Royal Holloway, University of London}

\begin{document}

\leveltitle{4F}{Difference and Joins, by Comparing Rows}
  {Programs track: follows FPIL}{After \levellink{3f}}
  {Next: \levellink{5f}}

\levelframe{4F}{Level 4F at a Glance}{%
  \code{dedup}, \code{diff} and \code{inter}; and \code{joinOn}, a join,
  as a selection of a product.}{%
  which sessions had nobody, and the size of a join, by computing.}{%
  write functions for every type whose values can be compared &
    \rds{fpil-classes} \\
  compare values with \code{==}, and test membership in a list &
    \rds{fpil-eqord} \\
}

% ---------------------------------------------------------------- theory
\theorypart{Equal rows, difference, and joins}

\begin{frame}
  \frametitle{Comparing Rows Needs an Equality Test}
  \defn{Definition 4F.1 (Equality test)}{An \emph{equality test} on $D$
    is a function $\mathit{eq} : D \times D \to \B$ with
    $\mathit{eq}(t, s) = \true$ exactly when $t = s$.}
  \defn{Definition 4F.2 (Membership, duplicates)}{With $\mathit{eq}$:
    \[
      \begin{array}{@{}l@{\;=\;}l@{\qquad}l@{}}
        \mathit{contains}\ [\,]\ t & \false \\
        \mathit{contains}\ (s :: l)\ t & \mathit{eq}(s, t) \text{ or }
          \mathit{contains}\ l\ t \\
        \mathit{dedup}\ [\,] & [\,] \\
        \mathit{dedup}\ (t :: l) & \mathit{dedup}\ l & \text{if }
          \mathit{contains}\ l\ t = \true \\
        \mathit{dedup}\ (t :: l) & t :: \mathit{dedup}\ l & \text{if }
          \mathit{contains}\ l\ t = \false \\
      \end{array}
    \]}
  So $\mathit{contains}\ l\ t = \true$ exactly when $t$ is in the table
  $l$ represents. In Lean, \code{==} is the equality test, from
  \code{BEq}.
  \theorytag{4F}
\end{frame}

\begin{frame}
  \frametitle{Difference Removes the Rows of Another Table}
  \defn{Definition 4F.3 (Difference, intersection)}{For tables $R$ and
    $S$ of the same rows: $R - S = \{t \in R \mid t \notin S\}$ and
    $R \cap S = \{t \in R \mid t \in S\}$.}
  Which sessions had nobody? With one-column tables:
  \[
    \pi_1[\rel{Session}] - \pi_2[\rel{Attended}] = \{1, 2, 3, 4\} - \{1,
    2, 3\} = \{4\}.
  \]
  On lists, a difference is a selection (Level 2F), with the condition
  ``not in $m$'': \
  $\mathit{diff}\ l\ m = \mathit{select}\ (t \mapsto \text{not }
  \mathit{contains}\ m\ t)\ l$. \ In SQL: \code{EXCEPT}.
  \defn{Fact 4F.4}{$R \cap S = R - (R - S)$. Level 3T proves it.}
  \slideref{\lnk{ahv5}}
  \theorytag{4F}
\end{frame}

\begin{frame}
  \frametitle{A Join Is a Selection of a Product}
  \defn{Definition 4F.5 (Join)}{For functions $k_1$ on the rows of $R$ and
    $k_2$ on the rows of $S$ (their \emph{keys}):
    \ $R \join_{k_1 = k_2} S = \{(t, s) \in R \times S \mid k_1(t) =
    k_2(s)\}$.}
  It is a selection of the product (Levels 2F and 3F), with the condition
  $\mathit{eq}(k_1(t), k_2(s))$. Joining on the week, with $k_1 = \pi_2$
  and $k_2 = \pi_1$:\par
  {\centering\small
    \smalltable{$\rel{Attended} \join_{\pi_2 = \pi_1}
      \rel{Session}$}{ll}{row of $\rel{Attended}$ & row of $\rel{Session}$}
      {(Ada, 1) & (1, Lists) \\ (Alan, 1) & (1, Lists) \\
       (Alan, 2) & (2, Proofs) \\ (Grace, 3) & (3, Types)}\par}
  \theorytag{4F}
\end{frame}

% ---------------------------------------------------------------- Lean
\leanpart{Build dedup, diff, inter and joinOn}

\begin{frame*}
  \frametitle{Your Turn: Build dedup}
  \leanblock{dedup.lean}
  \begin{itemize}
    \item \code{[BEq α]}: rows of type \code{α} have an equality test,
      \code{==}, and \code{l.contains row} uses it.
    \item \ex{4F.1} Complete \code{dedup} with Definition 4F.2, and prove
      the example.
  \end{itemize}
  \slideref{\lnk{fpil-classes} \rfsep \lnk{fpil-eqord}}
  \leantag{4F}
\end{frame*}

\begin{frame*}
  \frametitle{Your Turn: Build diff and inter}
  \leanblock{diff.lean}
  \begin{itemize}
    \item \ex{4F.2} \code{diff}: a selection, with \code{s.contains row}.
    \item \ex{4F.3} \code{inter}: from \code{diff} alone (Fact 4F.4). Then
      prove the example.
  \end{itemize}
  \leantag{4F}
\end{frame*}

\begin{frame*}
  \frametitle{Your Turn: Derive a Join from select and product}
  \leanblock{join.lean}
  \begin{itemize}
    \item \code{k₁} and \code{k₂} are the keys of Definition 4F.5: here,
      each row's week.
    \item \ex{4F.4} Complete \code{joinOn}: \code{select} the pairs of
      \code{product r s} whose keys are equal. Then prove the example.
  \end{itemize}
  \leantag{4F}
\end{frame*}

\end{document}
